\documentclass[11pt]{article}
\usepackage[margin=1in]{geometry}
\usepackage{amsmath,amssymb,amsthm,hyperref}
\newtheorem{theorem}{Theorem}
\newcommand{\Ree}{\operatorname{Re}}
\newcommand{\Imm}{\operatorname{Im}}
\title{High-frequency roots and diffusive boundary feedback}
\author{Principia Math In Progress}
\date{28 September 2026}
\begin{document}
\maketitle
\begin{abstract}
For the scalar transport--diffusion feedback system of Bastin, Coron
and Hayat, we prove that the spectral abscissa is nonnegative for every
positive diffusion coefficient. The characteristic roots approach the
odd imaginary multiples of $\pi$ at high frequency. This proves their
Conjecture 2, with a stronger quantifier in the diffusion coefficient.
\end{abstract}

\section*{Task 1: the conjecture and a stronger assertion}
The system in Section 6 of \cite{BCH} is
\[
 y_t+y_x-\eta y_{xx}=0,\qquad \widehat y_t+\widehat y_x=0,
 \qquad 0<x<1,
\]
with $\eta>0$ and boundary conditions
\[
 y(t,0)=y(t,1)-\widehat y(t,1),\qquad
 \widehat y(t,0)=y(t,0),\qquad y_x(t,1)=0.
\]
Write
\[
 z(s)=\sqrt{1+4\eta s},\qquad
 \lambda_1(s)=\frac{1+z(s)}{2\eta},\quad
 \lambda_2(s)=\frac{1-z(s)}{2\eta},
\]
using the principal square root. The source's characteristic function is
\[
 F_\eta(s)=
 \frac{(\lambda_1e^{\lambda_1}-\lambda_2e^{\lambda_2})(1+e^{-s})}
      {(\lambda_1-\lambda_2)e^{\lambda_1+\lambda_2}}.
\]
Set $\mathcal S_\eta=\{s:F_\eta(s)=1\}$ and
$\sigma_\eta=\sup\{\Ree s:s\in\mathcal S_\eta\}$.
Conjecture 2 asserts that for every $\varepsilon>0$ there is $\eta_1>0$
such that $\sigma_\eta>-\varepsilon$ whenever $0<\eta<\eta_1$.

\begin{theorem}\label{thm}
For every fixed $\eta>0$ there are characteristic roots
\[
 s_k=i(2k+1)\pi+o(1)\qquad(k\longrightarrow\infty).
\]
Consequently $\sigma_\eta\geq0$ for every $\eta>0$, and Conjecture 2 holds.
\end{theorem}

\section*{The transfer factor on high-frequency disks}
Factor $F_\eta(s)=H_\eta(s)(1+e^{-s})$, where
\begin{equation}\label{H}
 H_\eta(s)=\frac{\lambda_1 e^{-\lambda_2}-\lambda_2e^{-\lambda_1}}
                  {\lambda_1-\lambda_2}
 =\frac{1+z}{2z}e^{(z-1)/(2\eta)}
   -\frac{1-z}{2z}e^{-(z+1)/(2\eta)}.
\end{equation}
Fix $0<r<1/2$, let $\omega_k=(2k+1)\pi$, and consider
$s=i\omega_k+w$ with $|w|\leq r$. For large $k$,
$\Imm(1+4\eta s)\geq4\eta(\omega_k-r)>0$, so the principal square
root and $H_\eta$ are holomorphic on a neighborhood of this closed disk.
For a number $t$ in the upper half-plane,
\[
 \Ree\sqrt t=\sqrt{\frac{|t|+\Ree t}{2}}.
\]
It follows throughout the disk that
\begin{equation}\label{zbound}
 \Ree z\geq L_k:=\sqrt{2\eta\omega_k+\tfrac12-4\eta r},
 \qquad |z|^2\geq4\eta(\omega_k-r).
\end{equation}
In particular $L_k\to\infty$ and $|z|\geq2$ for all sufficiently large
$k$. For these $k$,
\[
 \left|\frac{1+z}{2z}\right|\geq\frac14,
 \qquad \left|\frac{1-z}{2z}\right|\leq\frac34.
\]
The reverse triangle inequality in \eqref{H} therefore gives the uniform
lower bound
\begin{equation}\label{B}
 |H_\eta(s)|\geq B_k:={1\over4}e^{(L_k-1)/(2\eta)}
                         -{3\over4}e^{-(L_k+1)/(2\eta)}
 \longrightarrow\infty.
\end{equation}
Thus $H_\eta$ has no zeros on these disks for large $k$, and its reciprocal
converges uniformly to zero there. All limits in this paragraph keep
$\eta$ fixed; no uniformity as $\eta\to0$ is needed.

\section*{Root existence and the spectral abscissa}
Because $e^{-i\omega_k}=-1$, the equation $F_\eta(s)=1$ in the disk is
equivalent to
\begin{equation}\label{G}
 G_k(w):=1-e^{-w}-H_\eta(i\omega_k+w)^{-1}=0.
\end{equation}
On $|w|=r$ the exponential series and the geometric series give
\begin{equation}\label{m}
 |1-e^{-w}|\geq r-\sum_{j=2}^{\infty}\frac{r^j}{j!}
 \geq r-\frac{r^2}{1-r}
 =:m_r=\frac{r(1-2r)}{1-r}>0.
\end{equation}
For all sufficiently large $k$, \eqref{B} implies
$|H_\eta(i\omega_k+w)^{-1}|<m_r$ on that circle.
Rouch\'e's theorem applied to \eqref{G} shows that $G_k$ and
$1-e^{-w}$ have the same number of zeros inside it. Since $r<1/2$,
the latter has exactly one zero, the simple zero $w=0$.
There is therefore exactly one characteristic root in
$|s-i\omega_k|<r$, counted with multiplicity.

First choose $r=1/4$ to define $s_k$ uniquely for all sufficiently large
$k$. Repeating the argument for any smaller positive radius shows that
these same roots satisfy $s_k-i\omega_k\to0$. In particular
$\Ree s_k\to0$, so $\sigma_\eta\geq0$. Since $\eta>0$ was arbitrary,
this proves Theorem \ref{thm}. For any $\varepsilon>0$ the desired
strict inequality follows from $\sigma_\eta\geq0>-\varepsilon$;
one may take, for example, $\eta_1=1$.

These roots also give actual separated solutions of the boundary system,
without a possible artifact from multiplying a characteristic equation.
At any of the roots, put
\[
 \phi(x)=\frac{-\lambda_2e^{\lambda_1(x-1)}
                         +\lambda_1e^{\lambda_2(x-1)}}{\lambda_1-\lambda_2},
 \qquad \psi(x)=H_\eta(s)e^{-sx}.
\]
The identities $\eta\lambda_j^2-\lambda_j-s=0$ give the differential
equations for $(y,\widehat y)=e^{st}(\phi,\psi)$. Moreover,
$\phi(1)=1$, $\phi'(1)=0$, $\phi(0)=\psi(0)=H_\eta(s)$, and
$\phi(0)=\phi(1)-\psi(1)$ follows from $H_\eta(s)(1+e^{-s})=1$.
The mode is nonzero since $\phi(1)=1$.

\section*{Exact reproducible disk certificates}
For rational $\eta>0$ and $0<r<1/2$, choose an integer $A>0$ and an
integer $N\geq0$ such that
\[
 P_N(A):=\sum_{j=0}^N\frac{A^j}{j!}>\frac4{m_r}+3.
\]
Using $\pi>3$, the following rational inequalities suffice for the
Rouch\'e comparison:
\[
 4\eta\bigl(3(2k+1)-r\bigr)>4,\qquad
 6\eta(2k+1)+\tfrac12-4\eta r>(1+2\eta A)^2.
\]
Indeed, they imply $|z|>2$, $L_k>1+2\eta A$ and
$B_k> P_N(A)/4-3/4>1/m_r$. The accompanying verifier checks
24 such certificates exactly, as well as the algebraic factorization
and mode identities. These are certified disks, not numerical root
approximations. The analytic argument proves the full parameter range.

\section*{Proposed follow-up tasks}
\begin{description}
\item[Task 2.] Determine whether every $\eta>0$ produces infinitely many
roots with strictly positive real part, rather than merely real parts
approaching zero.
\item[Task 3.] Derive a two-term asymptotic expansion for
$s_k-i(2k+1)\pi$, with an explicit remainder for fixed $\eta$.
\item[Task 4.] Replace the observer transport equation by
$\widehat y_t+c\widehat y_x=0$, $c>0$, and determine the corresponding
high-frequency root lattice and spectral-abscissa obstruction.
\item[Task 5.] For feedback $y(t,0)=a y(t,1)-b\widehat y(t,1)$ with
real $a,b$, classify which parameters retain roots approaching the
imaginary axis for every $\eta>0$.
\end{description}

\begin{thebibliography}{9}
\bibitem{BCH} G. Bastin, J.-M. Coron and A. Hayat,
\emph{Diffusion and robustness of boundary feedback stabilization of
hyperbolic systems}, Mathematics of Control, Signals, and Systems
\textbf{35} (2023), 159--185,
\href{https://doi.org/10.1007/s00498-022-00335-0}{doi:10.1007/s00498-022-00335-0}.
Conjecture 2 and equations (97)--(100), p.~179;
\href{https://arxiv.org/abs/2212.04879}{arXiv:2212.04879v1}, pp.~16--17.
\end{thebibliography}
\end{document}
